% =====================================================================
%  19.3  Searching the Tree                        sec:c13-search
%  Exports: thm:c13-descend  lst:c13-descend  lst:c13-kthseg
%           lst:c13-1179-soln  fig:c13-descend  prob:c13-dining
%  Solved: CF 1179C Serge and Dining Room (2200), live-verified.
%    Reduction: for a value v a dish of price v survives the queue iff the
%    suffix-sum S(v)=sum_{u>=v}(#dishes(u) - #pupils(u)) > 0; the answer is
%    the LARGEST such v, else -1.  Segment tree over values 1..1e6 storing
%    (sum, best=max suffix) and a DESCEND to the largest positive suffix.
%    Samples (verbatim from CF):
%      1) n=1 m=1 a=[1] b=[1] q=1 "1 1 100" -> 100
%      2) n=1 m=1 a=[1] b=[1] q=1 "2 1 100" -> -1
%      3) 4 6 / a=1 8 2 4 / b=3 3 6 1 5 2 / q=3 (1 1 1),(2 5 10),(1 1 6)
%         -> 8 / -1 / 4   (hand-traced: max-positive-suffix value)
%  Running numbers agree with 00_opener (fig:c13-tree).
% =====================================================================
\section{Searching the Tree}\label{sec:c13-search}

\begin{previouslybox}[Ch.~18 found answers by lifting; here we descend]
Chapter~18 closed its order-statistics section (\(\secref{c14-order}\)) with a
question ``which position is the \(k\)-th set bit?'' answered by \emph{lifting}
up the Fenwick tree -- a run of \(\text{lowbit}\)-sized hops that accumulates a
prefix until it is about to overshoot.  The very same search lives inside a
segment tree, but walks it \emph{downward}: start at the root, decide at every
node which child holds the answer, and recurse into just that one child.  No
range query, no canonical cover -- the tree is read as a decision procedure.
This section makes that exact move three times (first prefix \(\ge\) target,
\(k\)-th by counts, first element \(>x\)) and then spends it on a 2200-rated
CF problem whose answer is a threshold on a suffix sum.
\end{previouslybox}

\begin{topicintro}
A segment tree stores, at each node, an aggregate of its interval.  \emph{Query}
folds \(O(\log n)\) of those aggregates to answer a question about a
\emph{given} range.  \term{Descending} the tree asks the opposite kind of
question: \emph{which} range is the answer -- ``the first position whose prefix
sum reaches \(T\),'' ``the value of the \(k\)-th element,'' ``the first element
in \([l,r]\) greater than \(x\).''  Each is phrased as a predicate \(P(\mathit{node})\)
computable from a node's stored aggregate, decided top-down: at node \(v\) with
children \(L,R\), test whether the answer lies in \(L\) using only \(L\)'s
aggregate; if not, spend \(L\)'s aggregate to re-target and go to \(R\).  Because
exactly one child is entered per level, the walk is \(\bigO{\log n}\) -- the same
cost as a query, but it \emph{finds} the boundary instead of summing a range the
caller already knew.

\textbf{Prerequisites.} \(\secref{c13-core}\) (heap layout, node aggregates).
\(\secref{c13-combine}\) (a node may store a \emph{record}, which is exactly what
a descend predicate reads).  Ch.~18's \(\secref{c14-order}\) lifting as the
contrasting search.  Ch.~1's binary search for the mental model: a descend is
binary search where the ``midpoint test'' is a node aggregate in \(O(1)\).

\textbf{After this section you can}
\begin{itemize}[nosep]
  \item write \code{first\_prefix\_ge} and \code{kth} descends from memory and
        state the monotonicity each needs; \hfill\lv{3}--\lv{4}
  \item re-solve Ch.~18's \(k\)-th problem as a tree descend and compare the two
        walks line by line; \hfill\lv{4}
  \item design a node record whose aggregate \emph{drives} a top-down search
        (sum+best-suffix, count, max); \hfill\lv{4}--\lv{5}
  \item solve CF~1179C end to end -- a point update and a threshold answer in
        \(\bigO{\log V}\) each. \hfill\lv{5}
\end{itemize}
\end{topicintro}

\begin{recognitionbox}[Pattern card -- ``which position, found by walking down'']
\reckey{``the smallest \(i\) with prefix sum \(\ge T\)'' (needs monotone
prefixes); ``the \(k\)-th smallest / \(k\)-th active element''; ``the first
element of \([l,r]\) exceeding \(x\)''; ``the largest threshold \(v\) for which
some suffix aggregate stays positive''; any answer that is a \emph{position or
value to be located}, not a fixed range to be folded -- especially paired with
point updates so a prefix array cannot be reused.}
\recuse{
  \item a threshold on a running sum/count/max, found top-down -- this section;
  \item the predicate is monotone along the index (once true stays true), so a
        single-child descent is well-defined;
  \item the answer must reflect interleaved point updates -- the tree is the
        dynamic analogue of Ch.~1's static binary search.
}
\recnotuse{
  \item you already know the range and only need its fold -- that is a plain
        \(\coderef{c13-query}\), \(\secref{c13-core}\); no descent;
  \item the predicate is \emph{not} monotone in the index (a ``first local
        maximum'' style test) -- descend cannot commit to one child; reframe or
        binary-search-on-the-answer instead;
  \item no updates, pure static \(k\)-th in a range -- Ch.~21's merge-sort or
        persistent tree (\(\chref{c19}\)) answers in \(O(\log n)\) per query
        with \(O(1)\) amortised build effort.
}
One-line test: \emph{the answer is an index or value you must locate under a
monotone predicate, and the array changes -- descend the tree, \(\log n\) steps.}
\end{recognitionbox}

% ---------------------------------------------------------------
\subsection{The Descend Template}
% ---------------------------------------------------------------

Every search below is one shape.  A node \(v\) owns \([l,r]\), splits at
\(mid\), and stores an aggregate that tells us \emph{whether the target lies in
the left child}.  If it does we recurse left with the target unchanged; if not
we recurse right having \emph{consumed} the left child's contribution.

\begin{lstlisting}[caption={first position whose running prefix sum reaches target: a left-vs-right commit}, label={lst:c13-descend}]
// SUM tree: sum[v] = A[l]+..+A[r].  Requires all A[i] >= 0 so that the
// prefix sums are non-decreasing (else "first >= target" is ill-defined).
int first_prefix_ge(int v, int l, int r, long long target){
    if (l == r) return l;                        // accumulated >= target here
    int mid = (l + r) / 2;
    if (sum[2*v] >= target)                      // answer inside the left half
        return first_prefix_ge(2*v,   l,   mid,   target);
    return first_prefix_ge(2*v+1, mid+1, r, target - sum[2*v]); // consume left
}
// call first_prefix_ge(1, 1, n, T); O(log n): one child entered per level
\end{lstlisting}

\begin{lstlisting}[caption={k-th active element by counts -- Ch.18 lifting, re-done as a descend}, label={lst:c13-kthseg}]
// COUNT tree: cnt[v] = # active positions in [l,r].  Positions sorted by value.
int kth(int v, int l, int r, long long k){       // 1-based; needs k <= cnt[1]
    if (l == r) return l;                        // the l-th value is the k-th active
    int mid = (l + r) / 2;
    if (cnt[2*v] >= k)                           // k-th already in the left block
        return kth(2*v,   l,   mid,   k);
    return kth(2*v+1, mid+1, r, k - cnt[2*v]);   // skip the left block's count
}
\end{lstlisting}

\begin{theorem}{A monotone predicate descends in \(\bigO{\log n}\)}{c13-descend}
Suppose a node's stored aggregate suffices to decide, for any target, whether
the sought position lies in its left or right child (and, when it lies right, to
recompute the target from the left child's aggregate).  If the predicate
``the answer is \(\le i\)'' is monotone in \(i\), then a single top-down walk of
\(\le\lceil\log_2 n\rceil{+}1\) steps, entering one child per level, returns the
answer.  Its cost equals that of a range query but uses none of its branching:
there is exactly one recursive call at each node.
\end{theorem}
\begin{proofbox}
Monotonicity guarantees the target region is a contiguous prefix (or suffix), so
the left-vs-right test never needs to search both children: at \(v\), the left
child's aggregate is a faithful summary of every position \(\le mid\) (by
\defref{c13-node}'s contract -- a parent, hence the answer's predicate, is a
function of child aggregates).  Entering only the child that can contain the
boundary preserves the invariant ``the answer is inside the current node'' while
shrinking the interval by half; after \(\le\lceil\log_2 n\rceil{+}1\) halvings the
interval is a single leaf, which \emph{is} the answer.  The re-target
(\(\mathit{target}-\text{left aggregate}\)) keeps the predicate correct in the
child because the child only knows about its own positions.
\end{proofbox}

\noindent\textbf{In plain words.}  A range query is binary search's ``sum this
interval''; a descend is binary search's ``which interval.''  Ch.~18's lifting
does the same job from the \emph{leaf upward} by \(\text{lowbit}\) jumps; the tree
descend does it \emph{root downward} by midpoint choices.  They agree because
both are reading the same dyadic partition -- one from the array's binary
indices, one from the tree's shape.

\begin{figure}[H]
\centering
\begin{tikzpicture}[diagram, scale=0.9, transform shape]
  \node[boxnode, primaryblue, minimum size=8mm, font=\scriptsize] (r)  at (0,3)   {\([1..8]\)};
  \node[boxnode, difforange,  minimum size=8mm, font=\scriptsize] (la) at (-2.4,1.6){\([1..4]\)};
  \node[boxnode, successgreen,minimum size=8mm, font=\scriptsize] (rb) at (2.4,1.6) {\([5..8]\)};
  \node[boxnode, gray, minimum size=7mm, font=\scriptsize] (la1) at (-3.6,0.2){\([1..2]\)};
  \node[boxnode, gray, minimum size=7mm, font=\scriptsize] (la2) at (-1.2,0.2){\([3..4]\)};
  \node[boxnode, successgreen, minimum size=7mm, font=\scriptsize] (rb1) at (1.2,0.2){\([5..6]\)};
  \node[boxnode, successgreen, minimum size=7mm, font=\scriptsize] (rb2) at (3.6,0.2){\([7..8]\)};
  \draw[flow] (r)--(la); \draw[flow] (r)--(rb);
  \draw[flow, gray!40] (la)--(la1); \draw[flow, gray!40] (la)--(la2);
  \draw[flow] (rb)--(rb1); \draw[flow] (rb)--(rb2);
  \node[annot, font=\scriptsize, text=difforange, anchor=east] at (-2.4,2.4) {rejected};
  \node[annot, font=\scriptsize, text=successgreen, anchor=west] at (2.4,2.4) {taken};
  \node[annot, align=center, font=\scriptsize] at (0,-1.1)
    {left aggregate \(<\) target \(\Rightarrow\) go right,\\ consume the left, re-target by \(T-\mathit{sum}[\mathit{left}]\)};
\end{tikzpicture}
\caption{One step of a descend.  At node \([1..8]\) the left child \([1..4]\)'s
aggregate fails the predicate (its running count/sum does not reach the target),
so the whole left subtree is discarded \emph{at once} and we commit to
\([5..8]\), subtracting the left's contribution.  Three such commits resolve an
eight-leaf tree to a leaf -- that is \(\log_2 8=3\) decisions.}
\label{fig:c13-descend}
\end{figure}

% ---------------------------------------------------------------
\subsection{Three Descends, One Idea}
% ---------------------------------------------------------------

\term{First prefix \(\ge\) target.} \(\coderef{c13-descend}\), on a sum tree, all
values \(\ge0\).  This is the dynamic binary search over a prefix array -- Ch.~5's
\(\propref{c24-prefix}\) idea rebuilt to survive point updates.

\term{\(k\)-th by counts.} \(\coderef{c13-kthseg}\), on a tree of counts: the same
code Ch.~18 wrote as a lowbit lift, now a midpoint descend.  Feed it a multiset
(``insert/delete a value; what is the \(k\)-th smallest?'') by storing at each leaf
whether/how-many of that value are present.

\term{First element \(>x\) in \([l,r]\).} Descend a \(\max\)-tree: at \(v\), if
\(\mathit{mx}[\text{left}]\le x\) the left holds no candidate so go right; else go
left, and prune any node fully outside \([l,r]\) or with \(\mathit{mx}\le x\).  It
returns the leftmost position of \([l,r]\) whose value exceeds \(x\) in
\(\bigO{\log n}\) -- the engine behind ``next greater in a range'' sweeps.

% ---------------------------------------------------------------
\subsection{Problem Lab: CF 1179C Serge and Dining Room}
% ---------------------------------------------------------------

\solvedproblem{c13-dining}{Serge and Dining Room}{CF 1179C}
  {https://codeforces.com/problemset/problem/1179/C}{CF 2200}
  {max-suffix segment tree over values + descend to the largest positive suffix; point updates on dish and pupil counts}

\begin{problemstatement}[CF 1179C Serge and Dining Room]
There are \(n\) dishes with prices \(a_1..a_n\) and \(m\) pupils in a queue with
budgets \(b_1..b_m\) (values \(\le 10^6\)).  Each pupil, in queue order, buys the
\emph{most expensive} dish he can afford (each dish has one copy; once bought it
is gone); if he can afford none he leaves.  Serge, who can afford anything, buys
the \emph{most expensive remaining} dish after the whole queue has been served,
or gets nothing if no dish remains.  Process \(q\) point updates:
\texttt{1 i x} set dish \(a_i:=x\); \texttt{2 i x} set pupil \(b_i:=x\).  After
\emph{each} update, print the price of the dish Serge would buy, or \(-1\).

\textbf{Input.} \(n,m\le 3\cdot10^5\); \(1\le a_i,b_i,x\le 10^6\); \(q\le 3\cdot10^5\).

\textbf{Output.} One line per update.

\textbf{Examples.}
\begin{verbatim}
Input:               Output:      Input:          Output:   Input:
1 1                  100          1 1             -1        4 6
1                                 1                        1 8 2 4
1                                 1                        3 3 6 1 5 2
1 1 100                          2 1 100                  3
                                                             1 1 1
                                                             2 5 10
                                                             1 1 6
Output (3rd):
8
-1
4
\end{verbatim}
\end{problemstatement}

\begin{attemptgate}[Attempt it yourself first -- \lv{5}, give yourself one sitting]
Simulating the queue per update is \(\bigO{(n+m)\log}\) \emph{times} \(q\) -- far
too slow, and the queue order makes a naive multiset greedy miss the point.
Reframe away from \emph{which pupil takes which dish} and onto \emph{at each
price level, do dishes outnumber the pupils who can afford that price or more?}
\peek{a dish priced \(v\) is grabbed only by pupils with budget \(\ge v\); so
look at the values from high to low.}
\nudge{for a value \(v\), let \(D_{\ge}(v)\) be the number of dishes of price
\(\ge v\) and \(P_{\ge}(v)\) the pupils with budget \(\ge v\).  The number of
dishes priced \(\ge v\) left for Serge is \(S(v)=D_{\ge}(v)-P_{\ge}(v)\).  Serge
gets the \textbf{largest \(v\) with \(S(v)>0\).}  \(S\) is a \emph{suffix sum}.}
\fullhint{Build a tree over \emph{values} \(1..V\) (\(V=10^6\)) where leaf \(u\)
holds \(\mathit{delta}[u]=\#\text{dishes}(u)-\#\text{pupils}(u)\).  Each node
stores \(\mathit{sum}=\sum\mathit{delta}\) and \(\mathit{best}=\max\) suffix-sum.
An update moves one dish/pupil: adjust \(\mathit{delta}\) at two values (point
updates).  The answer is the \emph{largest value with positive suffix-sum}, found
by a \(\bigO{\log V}\) \textbf{descend}: prefer the right child (larger values)
when its \(\mathit{best}>0\), else go left carrying the right child's sum.  If the
root's \(\mathit{best}\le0\) the answer is \(-1\).  Full program below.}
\end{attemptgate}

\solstep{Understanding}
Fix a price \(v\).  Consider the dishes of price \(\ge v\): a pupil takes one of
them exactly when his budget is \(\ge v\) (he can afford such a dish, and greedily
prefers expensive).  So the count of price-\(\ge v\) dishes \emph{still on the
shelf} once everyone is served is \(S(v)=D_{\ge}(v)-P_{\ge}(v)\), independent of
the exact pairing.  Serge takes the most expensive survivor, i.e.\ the largest
\(v\) that still has one left: the largest \(v\) with \(S(v)\ge 1\) \emph{and} a
dish at exactly \(v\).  Because \(S(v)-S(v{+}1)=\mathit{delta}[v]\) and a
strict-positive suffix already forces a dish at the boundary \(v\), the answer is
precisely \(\max\{v: S(v)>0\}\).  Now \(S(v)\) is the suffix sum of
\(\mathit{delta}\) at \(v\), and \(\max\) over suffix sums is a node record -- this
section's descend, driven by \((\mathit{sum},\mathit{best})\).

\solstep{Key Observations}
\begin{observation}
The queue \emph{order} never matters.  Only the \emph{multiset} of budgets and of
prices does -- greedy ``most expensive affordable'' makes the count argument
\(S(v)=D_{\ge}(v)-P_{\ge}(v)\) exact.  This is why a value-indexed tree, not a
position-indexed one, is the right object.
\end{observation}
\begin{observation}
\(\mathit{best}\) = \emph{maximum suffix sum} is the record that makes a search
possible: a plain sum cannot tell you \emph{where} a suffix first goes positive.
The pair \((\mathit{sum},\mathit{best})\) is the minimal closing record for a
``largest positive suffix'' query -- a \(\secref{c13-combine}\) idea in service of
a \(\secref{c13-search}\) search.
\end{observation}
\begin{observation}
The update is \emph{two} point writes: a dish moving \(a_i:\)old\(\to x\) does
\(\mathit{delta}[\text{old}]-{=}1,\ \mathit{delta}[x]+{=}1\); a pupil does the
opposite sign.  Both are \(\bigO{\log V}\) re-folds, then one \(\bigO{\log V}\)
descend for the answer.
\end{observation}

\solstep{Think It Through With Me}
\begin{thinkbox}
\deadend{Simulate the queue with a sorted multiset of dishes each query}
{Each simulation is \(\bigO{(n+m)\log}\) and there are \(q\) of them --
\(9\cdot10^{10}\).  The multiset also re-does the whole greedy from scratch even
though only one dish or pupil changed.}
\deadend{Store a prefix-sum tree and binary-search \(S(v)>0\)}
{A prefix-sum (in the value order used here, a \emph{suffix}-sum) tree answers a
fold over a \emph{given} \(v\) in \(\log V\), but we need the largest \(v\)
satisfying a predicate -- binary-searching over value-space with a \(\log V\) fold
at each step is \(\bigO{\log^2 V}\).  The \((\mathit{sum},\mathit{best})\) descend
does it in a single \(\bigO{\log V}\) walk.}
\deadend{Use a Fenwick tree over the value counts}
{A BIT stores sums and can lift to a \(k\)-th, but ``largest positive suffix
sum'' needs the max-suffix \(\mathit{best}\) record, which has no inverse --
the same no-inverse wall Ch.~18 hit.  The segment tree's extra field is exactly
the point.}
\aha{Index by \emph{value}, store \((\mathit{sum},\mathit{max\text{-}suffix})\),
and \emph{descend} to the largest value whose suffix-sum is positive: prefer the
bigger-value child when its \(\mathit{best}>0\), else fold the right child's sum
into the accumulator and go left.}
\end{thinkbox}

\solstep{The Approach Ladder}
\begin{center}\renewcommand{\arraystretch}{1.25}
\begin{tabular}{@{}c p{0.22\linewidth} p{0.31\linewidth} l c@{}}
\toprule
\textbf{\#} & \textbf{Approach} & \textbf{Idea} & \textbf{Time} & \textbf{Verdict}\\
\midrule
1 & \lv{3}\ simulate queue & sorted dishes, greedy per query & \(\bigO{q(n{+}m)\log}\) & \tle \\
2 & \lv{4}\ BIT of deltas \(+\) binary search & prefix fold per \(v\) & \(\bigO{q\log^2 V}\) & \tle-ish \\
3 & \lv{4}\ sum-only segtree & no max-suffix field & \(\bigO{q\log^2 V}\) & slow \\
4 & \lv{5}\ \textbf{\((\mathit{sum},\mathit{best})\) tree} & descend to largest positive suffix & \(\bigO{(n{+}q)\log V}\) & \ac \\
\bottomrule
\end{tabular}\end{center}

\solstep{Solution}
Tree over \(1..V\) (\(V=10^6{+}1\)).  Leaf \(u\): \(\mathit{sum}=\mathit{delta}[u]\),
\(\mathit{best}=\mathit{delta}[u]\).  Merge (left \(A\)=smaller values, right
\(B\)=larger): \(\mathit{sum}=A.\mathit{sum}+B.\mathit{sum}\);
\(\mathit{best}=\max(B.\mathit{best},\ A.\mathit{best}+B.\mathit{sum})\) -- a suffix
starts inside the right child, or uses all of the right plus a suffix of the left.
Answer via \(\text{descend}\): from the root with accumulator \(\mathit{acc}=0\),
if \(\mathit{best}[\text{right}]+\mathit{acc}>0\) go right (bigger values); else go
left with \(\mathit{acc}+\mathit{sum}[\text{right}]\).  Print the leaf value, or
\(-1\) when \(\mathit{best}[1]\le 0\).

\solstep{Proof of Correctness}
\begin{proofbox}
\textbf{Record invariant.} For node over \([l,r]\), \(\mathit{best}=\max_{v\in[l,r]}
\sum_{u=v}^{r}\mathit{delta}[u]\) (max suffix of the node) and
\(\mathit{sum}=\sum_{u=l}^{r}\mathit{delta}[u]\).  Leaves are trivial; the merge
cases cover ``suffix inside right child'' and ``all of right child plus a suffix of
left child,'' the only shapes a suffix of \([l,r]\) can take.  \textbf{Answer:}
\(S(v)=\sum_{u=v}^{V}\mathit{delta}[u]=\) (suffix from \(v\) at the root), so Serge's
price is \(\max\{v:S(v)>0\}\).  \(\coderef{c13-1179-soln}\)'s \text{descend}
maintains the invariant ``\(\mathit{best}[\mathit{node}]+\mathit{acc}>0\)'' (a
positive suffix exists in this node once everything to its right, summed as
\(\mathit{acc}\), is added) and at each step enters the unique child that can
contain the \emph{largest} such \(v\): right if a positive suffix already lies
there, else left with the right block's sum folded into \(\mathit{acc}\).  By
\thmref{c13-descend} this is correct and \(\bigO{\log V}\); a non-positive root
\(\mathit{best}\) means every \(S(v)\le0\), i.e.\ no survivor, giving \(-1\).
\textbf{Updates:} each dish/pupil change is \(\pm1\) on two \(\mathit{delta}\)
leaves, correctly re-folding their ancestors, so the record invariant holds for
the next query.
\end{proofbox}

\solstep{Implementation}
\begin{lstlisting}[caption={CF 1179C full solution: value-indexed (sum,max-suffix) tree + descend}, label={lst:c13-1179-soln}]
#include <bits/stdc++.h>
using namespace std;
const int V = 1000002;            // values 1..1e6 (+1 safety)
const int MX = 300005;
int n, m, q, a[MX], b[MX];
int cD[V], cP[V];                 // dishes/pupils present at each value
long long sum[4*V], best[4*V];

void pull(int v){
    sum[v]  = sum[2*v] + sum[2*v+1];
    best[v] = max(best[2*v+1], best[2*v] + sum[2*v+1]); // max suffix: right, or right+left-suffix
}
void build(int v, int l, int r){
    if (l == r) { sum[v] = cD[l] - cP[l]; best[v] = sum[v]; return; }
    int mid = (l + r) / 2;
    build(2*v, l, mid); build(2*v+1, mid+1, r); pull(v);
}
void add(int v, int l, int r, int pos, long long d){
    if (l == r) { sum[v] += d; best[v] = sum[v]; return; }
    int mid = (l + r) / 2;
    if (pos <= mid) add(2*v, l, mid, pos, d);
    else            add(2*v+1, mid+1, r, pos, d);
    pull(v);
}
int descend(int v, int l, int r, long long acc){   // best[v]+acc > 0 is guaranteed
    if (l == r) return l;
    int mid = (l + r) / 2;
    if (best[2*v+1] + acc > 0)                     // a positive suffix lives in the larger half
        return descend(2*v+1, mid+1, r, acc);
    return descend(2*v, l, mid, acc + sum[2*v+1]); // else fold the right block and go left
}
int main(){
    scanf("%d%d", &n, &m);
    for (int i=1;i<=n;i++){ scanf("%d",&a[i]); cD[a[i]]++; }
    for (int i=1;i<=m;i++){ scanf("%d",&b[i]); cP[b[i]]++; }
    build(1,1,V);
    scanf("%d",&q);
    while (q--){
        int ty,i,x; scanf("%d%d%d",&ty,&i,&x);
        if (ty==1){                                 // dish i : old a[i] -> x
            cD[a[i]]--; add(1,1,V,a[i],-1);
            cD[x]++;    add(1,1,V,x,+1);
            a[i]=x;
        } else {                                    // pupil i : old b[i] -> x
            cP[b[i]]--; add(1,1,V,b[i],+1);
            cP[x]++;    add(1,1,V,x,-1);
            b[i]=x;
        }
        if (best[1] <= 0) printf("-1\n");
        else printf("%d\n", descend(1,1,V,0));
    }
}
\end{lstlisting}

\solstep{Trace the Third Official Sample}
\begin{tracebox}[CF 1179C sample 3: dishes 1 8 2 4, pupils 3 3 6 1 5 2]
Initial \(\mathit{delta}[u]=\#\text{dishes}(u)-\#\text{pupils}(u)\):
\(\delta_1{=}0\), \(\delta_2{=}0\), \(\delta_3{=}{-}2\), \(\delta_4{=}{+}1\),
\(\delta_5{=}{-}1\), \(\delta_6{=}{-}1\), \(\delta_8{=}{+}1\), else \(0\).  Suffix
sums \(S(v)=\sum_{u\ge v}\delta_u\): \(S_8{=}1\), \(S_6{=}0\), \(S_5{=}{-}1\),
\(S_4{=}0\), \(S_3{=}{-}2\).
\textbf{Query \texttt{1 1 1}} (dish~1 stays \(1\)): deltas unchanged, largest
positive \(S\) is \(S_8=1\) \(\Rightarrow\) Serge gets \(8\).  Output \(8\).
\textbf{Query \texttt{2 5 10}} (pupil~5: \(5\to10\)): \(\delta_5{+}{=}1\to0\),
\(\delta_{10}{-}{=}1\to-1\).  Now \(S_{10}{=}{-}1\), \(S_8{=}1{-}1{=}0\) and no suffix
is positive -- the root \(\mathit{best}\le0\).  Output \(-1\).
\textbf{Query \texttt{1 1 6}} (dish~1: \(1\to6\)): \(\delta_1{-}{=}1\to-1\),
\(\delta_6{+}{=}1\to0\).  Recompute suffixes downward: \(S_8{=}0\) (the \(+1\) at
\(8\) is cancelled by \(\delta_{10}{=}{-}1\)), then \(S_7{=}S_6{=}S_5{=}0\), and
\(S_4{=}\delta_4{+}S_5{=}1{+}0{=}1{>}0\) while \(S_3{=}{-}2{+}1{=}{-}1\).  The
largest value with a positive suffix is \(v=4\).  Output \(4\).
Final outputs \texttt{8, -1, 4} -- the official answers.
\end{tracebox}

\solstep{Complexity}
\begin{complexitybox}
Build \(\bigO{V}\) with \(V=10^6{+}2\) (about \(4\cdot10^6\) tree slots,
\(\sim 64\) MB for two \code{long long} arrays -- inside \(256\) MB, but size the
arrays \(4V\), not \(2V\)).  Per update: two point writes \(+\) one descend,
each \(\bigO{\log V}\approx 20\); with \(q,n,m\le3\cdot10^5\) the total is
\(\sim 1.2\cdot10^7\) operations -- \textbf{AC} at the \(4\) s limit.  The
\(\mathit{delta}\) counts need only \(|{\cdot}|\le\max(n,m)\le3\cdot10^5\), so
\code{int} suffices for the counts but \code{long long} for \(\mathit{sum}\) is
harmless.
\end{complexitybox}

\begin{hardeningbox}
\begin{itemize}[nosep,leftmargin=1.4em]
  \item \textbf{Index by value, not position.}  The tree's leaves are the prices
        \(1..V\); \(\mathit{delta}\) counts dishes minus pupils at each price.
        A position-indexed tree would have to re-scan the queue.
  \item \textbf{Two point writes per update.}  Forgetting the \emph{old} value's
        decrement (or mishandling type~2's sign flip) is the classic WA.
  \item \textbf{The \(\mathit{best}\) merge order.}  A suffix is ``right child
        fully used, then a left suffix'': \(\mathit{best}=\max(B.\mathit{best},
        A.\mathit{best}+B.\mathit{sum})\).  Swapping \(A,B\) silently answers the
        wrong side.
  \item \textbf{Guard \(\mathit{best}[1]\le0\) before descending.}  The descend
        assumes a positive suffix exists; with none it returns a spurious value
        instead of \(-1\).
\end{itemize}
\end{hardeningbox}

% ---------------------------------------------------------------
\subsection{Common Pitfalls}
% ---------------------------------------------------------------

\begin{warningbox}[Trap 1: descending a non-monotone predicate]
\(\coderef{c13-descend}\) needs prefix sums monotone (all \(A[i]\ge0\)); with a
negative element a ``first prefix \(\ge T\)'' can be followed by a smaller prefix,
so ``commit to the left if it reaches \(T\)'' is no longer correct.  Check the
monotonicity \thmref{c13-descend} requires before reusing a descend template.
\end{warningbox}

\begin{warningbox}[Trap 2: a sum-only tree where a record is required]
Searching for ``largest positive suffix'' with a tree that stores only \(\mathit{sum}\)
forces a second binary search over the value axis (\(\bigO{\log^2 V}\)) or fails.
The \(\mathit{best}\) field is what buys the single \(\bigO{\log V}\) walk -- it is
\(\secref{c13-combine}\)'s lesson: \emph{decide the node record from the search you
need}.
\end{warningbox}

\begin{warningbox}[Trap 3: \(\bigO{\log^2}\) when \(\bigO{\log}\) is available]
Wrapping a fold-based predicate in an outer binary search (as in ladder rung~2/3)
is often \emph{correct} but a factor of \(\log\) too slow at \(3\cdot10^5\).
The whole value of a descend is folding the search into one root-to-leaf walk.
\end{warningbox}

\begin{warningbox}[Trap 4: sign errors on the two value updates]
Type~1 (dish) and type~2 (pupil) adjust \(\mathit{delta}\) with \emph{opposite}
signs (a dish adds \(+1\) to its new price, a pupil subtracts \(1\) from his new
budget's price because \(\mathit{delta}=\#\text{dishes}-\#\text{pupils}\)).  One
wrong sign and every subsequent answer drifts -- pass samples, fail tests.
\end{warningbox}

% ---------------------------------------------------------------
\subsection{Practice Problems}
% ---------------------------------------------------------------

\begin{stuckbox}
\begin{itemize}[nosep,leftmargin=1.4em]
  \item Is your predicate monotone along the value/axis you descend?
  \item Which record field (count, sum, max-suffix, max) drives the left-vs-right
        commit at each node?
  \item Did you enter \emph{one} child per level (a search) rather than both (a
        query)?
  \item For CF~1179C: are the two \(\mathit{delta}\) point updates' signs correct
        for the query type?
  \item Is the tree sized \(4V\), and is \(\mathit{best}[1]\le0\) checked first?
\end{itemize}
\end{stuckbox}

\begin{practiceset}
  \practice{Misha and Permutations Summation}{CF 501D}
    {https://codeforces.com/problemset/problem/501/D}
    {CF 2000 L4}
    {Add two permutations by their lexicographic rank modulo \(n!\) -- via the
     Lehmer / factorial code, whose every step is exactly a \emph{\(k\)-th by
     counts} search over the set of not-yet-used values (digit \(i\) = how many
     unused values are below \(p_i\); decode = the \(k\)-th unused value).
     Do it \emph{both} ways: Ch.~18's frequency-BIT lowbit lift
     (\(\secref{c14-order}\)) and this section's \(\coderef{c13-kthseg}\) tree
     descend, and watch them trace the same dyadic path from opposite ends.}
  \practice{Army Creation}{CF 813E}
    {https://codeforces.com/problemset/problem/813/E}
    {CF 2200 L5}
    {Given types \(a_1..a_n\), answer online-obfuscated ranges \([l,r]\): the
     largest army using \(\le k\) warriors of each type.  Reduce each type to a
     running count and locate, inside \([l,r]\), how many positions are among the
     first \(k\) copies of their type -- a threshold/``how far do I go before the
     \(k{+}1\)-th copy appears'' search.  A segment tree over positions keyed by
     that cut point, searched as in \(\secref{c13-search}\), answers each range in
     \(\bigO{\log n}\) (the full offline/persistent treatment is Ch.~21).}
\end{practiceset}

\paragraph{Back-cites and forwards.}
\(\secref{c13-core}\) supplies the aggregates a descend reads and
\(\secref{c13-combine}\) the records (here \((\mathit{sum},\mathit{best})\)) that
make a search possible.  Chapter~18's \(\secref{c14-order}\) is the twin search
done by lowbit lifting -- practice CF~501D asks for both.  Chapter~1's
\(\chref{c01}\) binary search is the static ancestor of \thmref{c13-descend}.
Forward: \(\secref{c13-iter}\) re-implements a descend on the bottom-up \(2n\)
tree (the commit test reads the same aggregates); \(\secref{c13-static}\) marks
where a static \(k\)-th moves to Ch.~21's merge-sort/persistent trees
(\(\chref{c19}\)).

% ---------------------------------------------------------------
\begin{recallbox}[Section 19.3 in 10 seconds]
A \textbf{descend} answers ``\emph{which} position'' instead of ``fold this
range'': at each node, read one child's aggregate to decide whether the answer is
left or right, commit to \emph{one} child, \(\bigO{\log n}\).  Requires a
\textbf{monotone} predicate (\thmref{c13-descend}) and the right node record --
sum for first-prefix-\(\ge\), counts for \(k\)-th, max for first-\(>x\).
Ch.~18 lifts (leaf up), this descends (root down): same search, opposite ends.
CF~1179C: index by value, store \((\mathit{sum},\mathit{max\text{-}suffix})\),
descend to the largest positive suffix; \(\mathit{delta}[v]=\#\text{dishes}(v)-
\#\text{pupils}(v)\), two point writes per update.
\selftest{
  \item Which three predicates do \coderef{c13-descend}, \coderef{c13-kthseg}, and
        the max-tree ``first \(>x\)'' each require to be monotone?
  \item State the \((\mathit{sum},\mathit{best})\) merge for a maximum-suffix record.
  \item Why is a sum-only tree not enough to answer ``largest positive suffix'' in \(\log n\)?
  \item In CF~1179C, why does a pupil change flip the sign of the \(\mathit{delta}\) writes?
  \item How do Ch.~18's lowbit lifting and this descend trace the same path?
}
\end{recallbox}
